% #############################################################################
% This is Chapter 1
% !TEX root = ../main.tex
% #############################################################################
% Change the Name of the Chapter i the following line
\fancychapter{Introduction}
% The following line allows to ref this chapter
\label{chap:intro}
% o maravilhoso mundo dos agentes
Agentic applications~\cite{agentic-survey}, software systems that combine one or more large language models (LLMs) with tools, memory, retrieval systems, and external services to autonomously plan and execute multi-step tasks, have rapidly evolved from research prototypes into production systems. They now power coding assistants (e.g. Devin~\cite{devin}, GitHub Copilot Workspace~\cite{copilot-workspace}, Cursor Agent Mode~\cite{cursor}), web and browser agents (e.g. WebArena~\cite{webarena}, BrowserUse~\cite{browser-use}, OpenDevin Web~\cite{opendevin}), enterprise automation platforms (e.g. Salesforce Einstein GPT Agent~\cite{salesforce-agentforce}, Microsoft AutoGen~\cite{autogen}, CrewAI~\cite{crewai}), customer-support workflows (e.g. Zendesk AI Agent~\cite{zendesk-ai-agent}, Intercom Fin~\cite{intercom-fin}, Freshdesk Freddy~\cite{freshdesk-freddy}, Ada~\cite{ada-support}), and security tooling (e.g. Takedown (vulnerability discovery)~\cite{takedown}, Reaper (security scanning)~\cite{reaper-scanner}, Garak (red-teaming)~\cite{garak2024}, Giskard~\cite{giskard}, PromptGuard~\cite{promptguard2025}). Standardized interfaces such as the Model Context Protocol (MCP)~\cite{mcp-blog} and growing connector ecosystems have further accelerated adoption by allowing agents to discover and invoke third-party capabilities with minimal integration effort. The result is a new class of software whose control flow is increasingly generated dynamically by LLMs at run time rather than explicitly encoded in source code.

% o mma é afinal perigoso
This flexibility creates security challenges that differ fundamentally from those of both conventional software and standalone LLMs. Because agent planning and execution are mediated through natural-language context, the distinction between data and executable instructions becomes blurred: content consumed by an agent, including retrieved documents, tool outputs, memory entries, and web content, may influence or redirect behavior. Recent incidents demonstrate the practical consequences of this design shift. Researchers have shown indirect prompt injection attacks that exfiltrate private repositories through MCP integrations~\cite{mcp-blog}, browser-agent compromises that hijack authenticated sessions~\cite{pleasefix}, and vulnerabilities in coding and enterprise agents that enable unauthorized execution or data leakage~\cite{takedown}. As agents acquire broader privileges and access to external systems, the consequences of these failures increase substantially.

% RW
Research and industry have responded rapidly to these risks. Emerging taxonomies identify attack vectors including prompt injection~\cite{greshake, pleasefix, agentdojo}, malicious data injection~\cite{agent-card-poisoning}, memory poisoning~\cite{memory-poison-sunil}, and tool abuse~\cite{tooltweak}, while frameworks such as the OWASP Top 10 for LLM Applications~\cite{owasp} and MITRE ATLAS~\cite{mitre-atlas} catalog adversarial techniques and associated risks. Recent work has also proposed defenses spanning guardrails~\cite{xiang2024guardagent,pro2guard}, information-flow control~\cite{costa2025fides}, runtime monitoring~\cite{pro2guard}, privilege separation~\cite{progent}, identity management~\cite{openid-identity}, and human-in-the-loop validation~\cite{hula}, while benchmarks such as AgentDojo~\cite{agentdojo} and Agent Security Bench~\cite{asb} measure attack effectiveness under controlled settings. However, an important question remains unanswered: how prevalent are these risk-bearing patterns and their corresponding defenses in real agentic software?



This paper addresses that gap through an empirical prevalence study of security-relevant design patterns in open-source agentic applications implemented in Python and JavaScript. We focus on these ecosystems because they host many widely used agent-development frameworks and open-source applications. We assemble the \underline{A}gent \underline{R}isks and \underline{D}efenses \underline{A}rchive (ARDA), a dataset of [X] applications collected from GitHub and selected to capture diversity across frameworks, application domains, and autonomy levels within these ecosystems. We operationalize existing agent-security taxonomies into concrete, statically detectable signals representing both risk exposure and defensive mechanisms, and implement a static-analysis toolkit that detects these patterns at scale. This enables us to measure not only how frequently different risk categories appear, but also whether applications deploying potentially dangerous capabilities adopt corresponding mitigations; for example, whether execution-capable agents are more likely to implement guardrails, privilege separation, or runtime controls.

Our analysis is intentionally structural rather than exploit-driven. Static analysis identifies potential exposure and defensive coverage rather than confirmed exploitability; this distinction is important because agentic vulnerabilities are behavioral, context dependent, and model sensitive. Our objective is therefore not to prove exploitability, but to characterize the security posture and defensive maturity of contemporary Python- and JavaScript-based agentic ecosystems at scale.

[Findings placeholder.]

Beyond the prevalence study itself, we believe the resulting corpus provides a foundation for future benchmark construction. By identifying and annotating risk and defense sites within real applications, \textsc{AgentVulnDB} could support benchmark suites for evaluating vulnerability-detection tools and defense verification systems, analogous to benchmark collections used in traditional software security. We leave benchmark construction and dynamic exploit validation to future work.

The remainder of the paper is organized as follows. Section~\ref{sec:overview} gives an overview of the principal risks in agentic applications and their typical defenses. Section 3 introduces our dataset. Section 4 describes our static-analysis toolkit for detecting potential risks and the corresponding defenses. Section 5 presents the results. Section 6 surveys related work, and Section 7 concludes and discusses directions for future research.